\documentclass[11pt]{article}

% ---------------------------------------------------------------------------
% Thesis-style description of the algorithms in "Gibbs sampler code Optimized"
% and the quantities measured by run_test.py.
% ---------------------------------------------------------------------------

\usepackage[utf8]{inputenc}
\usepackage[T1]{fontenc}
\usepackage{amsmath,amssymb,amsthm}
\usepackage{mathtools}
\usepackage{booktabs}
\usepackage{algorithm}
\usepackage{algpseudocode}
\usepackage{hyperref}
\usepackage{geometry}
\geometry{margin=1in}

\newcommand{\Tr}{\operatorname{Tr}}
\newcommand{\R}{\mathbb{R}}
\newcommand{\PSD}{\succeq 0}
\newcommand{\id}{\mathbb{I}}
\newcommand{\eps}{\varepsilon}
\newcommand{\dist}{D_{\mathrm{tr}}}
\DeclareMathOperator*{\argmax}{arg\,max}

\title{A Quantum SDP Primal Oracle Driven by Discrete-Time Quantum
Gibbs Sampling:\\ Implementation and Benchmark Design}
\author{}
\date{}

\begin{document}
\maketitle

\begin{abstract}
This note documents the algorithms in \texttt{Gibbs sampler code Optimized} and
the benchmark driver \texttt{run\_test.py}. The code couples the SDP
\emph{primal oracle} of van Apeldoorn and Gily\'en~\cite{vAG2019} with a
classical simulation of the discrete-time quantum Gibbs sampler of Gily\'en,
Chen, Doriguello and Kastoryano~\cite{GCDK2024}. At each oracle iteration a
thermal state of a Hamiltonian built from the SDP data is prepared by iterating
a CPTP channel; we measure how many channel applications are required to reach
trace-distance tolerance. We describe the SDPA conversion (including a CVXPY
dual solve for the trace bound~$R$), the oracle loop, the sampling channel, the
distinction between jump operators and the initial state, and the four benchmark
experiments.
\end{abstract}

\tableofcontents

% ===========================================================================
\section{Overview of the pipeline}
% ===========================================================================

\begin{enumerate}
  \item \textbf{Conversion} (\texttt{sdp\_conversion.py}): SDPLIB \texttt{.dat-s}
  files $\to$ oracle \texttt{.npz} instances. By default CVXPY solves the SDPA
  dual to set $R=\text{margin}\cdot\Tr(Y^\star)$ (Section~\ref{sec:Rcvxpy}).
  \item \textbf{Primal oracle} (\texttt{primal\_oracle\_quantum\_v1cube.py}):
  one MMW feasibility test at calibrated threshold $g$
  (Section~\ref{sec:oracle}).
  \item \textbf{Gibbs sampler} (\texttt{gibbs\_sampler\_quantum\_v2cube.py}):
  discrete-time channel toward $\rho\propto e^{-H}$ (Section~\ref{sec:gibbs}).
  \item \textbf{Benchmark} (\texttt{run\_test.py}): $g$-calibration and
  experiments 2.1.1--2.1.4 (Section~\ref{sec:measure}).
\end{enumerate}

Notation: $n$ = SDP matrix dimension; $M$ = number of oracle inequalities;
channel dimension $d=n$.

% ===========================================================================
\section{Problem reformulation}
\label{sec:conversion}
% ===========================================================================

\subsection{Source format (SDPA sparse)}
Symmetric block-diagonal $F_0,\ldots,F_m$, objective $c\in\R^m$:
\begin{align}
\text{(P)}\quad & \min\; c^\top x
  \quad\text{s.t.}\quad \textstyle\sum_i F_i x_i - F_0 = X \PSD,\\
\text{(D)}\quad & \max_{Y\PSD}\; \Tr(F_0 Y)
  \quad\text{s.t.}\quad \Tr(F_i Y)=c_i,\ i=1,\ldots,m .
\end{align}
Under strong duality, $\text{(P)}_{\mathrm{opt}}=\text{(D)}_{\mathrm{opt}}$; the
metadata CSV lists this common value as \textsc{opt}.

\subsection{Oracle form}
Feasibility version of
\begin{equation}
\label{eq:oracleform}
\max\; \Tr(CX)\quad\text{s.t.}\quad
\Tr(A_j X)\le b_j,\ A_1=\id,\ b_1=R,\quad X\PSD .
\end{equation}
Dual reduction ($X$ plays the role of $Y$): $C=F_0$; each $\Tr(F_iY)=c_i$
becomes $\Tr(F_iX)\le c_i$ and $\Tr(-F_iX)\le -c_i$; prepend $\Tr(X)\le R$.
Constraint stack $[\id,F_1,-F_1,\ldots]$ with $M=1+2m$. The objective enters
only inside the oracle as $A_0=-C$, $b_0=-g$.

\subsection{Trace bound $R$ from the dual optimizer}
\label{sec:Rcvxpy}

The trace inequality is \emph{not} part of the SDPA dual. It is added so the
oracle works on $X'=X/R$ with $\Tr(X')\le 1$. For the bound to be valid and as
tight as possible we need $R\ge \Tr(Y^\star)$ for an optimal dual $Y^\star$.

\paragraph{Default (CVXPY).}
\texttt{sdp\_conversion.py} solves~\eqref{eq:oracleform} \emph{without} the
trace constraint---i.e.\ the pure dual~(D)---using CVXPY:
\begin{equation}
\label{eq:dualsolve}
\max_{Y\PSD}\ \Tr(F_0 Y)\quad\text{s.t.}\quad \Tr(F_i Y)=c_i .
\end{equation}
From the solution $Y^\star$ (symmetrized numerically),
\begin{equation}
\label{eq:Rcvxpy}
R = R_{\mathrm{margin}}\cdot \Tr(Y^\star),
\qquad R_{\mathrm{margin}}=1.1\ \text{(default)},
\end{equation}
and $\textsc{opt}_{\mathrm{cvxpy}}=\Tr(F_0 Y^\star)$ is stored as the reference
optimum. Solvers are tried in order CLARABEL $\to$ SCS $\to$ ECOS. Conversion of
all 15 current instances takes on the order of minutes (one-time cost).

\paragraph{Why this is better than $R=2|\textsc{opt}|$.}
The old heuristic conflated the \emph{objective value} with the \emph{trace} of
the optimizer. There is no inequality $\Tr(Y^\star)\le 2|\textsc{opt}|$ in
general. Examples from the converted set:
\begin{itemize}
  \item \texttt{hinf4}: $2|\textsc{opt}|\approx 550$ but $\Tr(Y^\star)\approx
  66$, so $R$ drops from ${\sim}550$ to ${\sim}72$ and
  $\theta=\eps/(2R)$ grows ${\sim}8\times$.
  \item \texttt{ThetaPrimeER23}: $\Tr(Y^\star)$ is large (${\sim}7.5\times
  10^4$), so $R$ is \emph{larger} than the old heuristic (${\sim}193$): the
  certified ball is honest but $\theta$ remains small on those instances.
\end{itemize}
The CSV \textsc{opt} is retained as \texttt{opt\_csv} for cross-checking;
\texttt{opt} prefers the solver value when available.

\paragraph{Fallback and override.}
If CVXPY fails: $R = R_{\mathrm{heur}}\cdot\max(|\textsc{opt}_{\mathrm{csv}}|,1)$
with default $R_{\mathrm{heur}}=2$. \texttt{--R} forces a manual $R$ for all
instances. \texttt{--no-R-from-solver} disables~\eqref{eq:dualsolve}.

Stored fields per instance: \texttt{R}, \texttt{tr\_Y\_star}, \texttt{R\_source}
(\texttt{cvxpy} / \texttt{heuristic} / \texttt{override}), \texttt{opt},
\texttt{opt\_csv}.

% ===========================================================================
\section{The primal oracle}
\label{sec:oracle}
% ===========================================================================

\subsection{Feasibility and Gibbs Hamiltonian}
Threshold $g$ adds $A_0=-C$, $b_0=-g$. Dual weights $y\ge 0$ and
\begin{equation}
\theta=\frac{\eps}{2R},\qquad
M=\sum_k y_k\,\tilde A_k,\qquad
\rho_n=\frac{e^{-M}}{\Tr\,e^{-M}}.
\end{equation}
Paper schedule $T=\lceil\ln n/\theta^2\rceil$; the benchmark caps iterations via
\texttt{--target-iters} and \texttt{--max-oracle-iters}.

\subsection{$(n{+}1)$ embedding}
Padded $\tilde\rho=\mathrm{diag}(X',\omega)$ with $\omega=1-\Tr(X')$ is used in
the analysis, but block-diagonal jumps do not mix $\omega$. The channel runs on
$\rho_n$; then
$\omega=1/(Z_n+1)$, $X'=(1-\omega)\rho_n$ with stable $\log Z_n$.

\subsection{Violation selection}
$\theta$-feasible when $\max_k(\Tr(\tilde A_k X')-\tilde b_k)\le\theta$.
Otherwise $y_j\mathrel{+}=\theta$, $M\mathrel{+}=\theta\tilde A_j$, with
\texttt{max} (largest violation) or \texttt{random} (uniform over violated
indices).

\subsection{Gibbs modes}
\begin{description}
  \item[\texttt{exact}] $\rho_n=e^{-M}/\Tr\,e^{-M}$; used for $g$-calibration.
  \item[\texttt{mcmc}] iterate $\mathcal{M}$ until
  $\dist(\sigma,\rho_n)\le\theta$; benchmark cost.
\end{description}
Each oracle iteration restarts from $\sigma_0=\id/n$ (no warm start) so reported
step counts are from-scratch mixing times at that iteration's $M$.

% ===========================================================================
\section{The discrete-time quantum Gibbs sampler}
\label{sec:gibbs}
% ===========================================================================

\subsection{Target, jumps, and initial state (GCDK)}
\label{sec:jumps-vs-init}

Following~\cite{GCDK2024} (Theorem~1, Eq.~(5); coherent accept, Eq.~(11)):

\begin{itemize}
  \item \textbf{Hamiltonian $H=M$} sets the \emph{fixed point}
  $\rho=e^{-\beta H}/\Tr(e^{-\beta H})$, $\beta=1$.
  \item \textbf{Jump operators} $\{A_a\}$ (here: SDP objective $C$ and constraint
  generators $F_i$, with $\id$ and redundant $\pm F_i$ dropped) define
  \emph{which directions} in state space the channel can move. Each $A_a$ is
  Bohr-reweighed in the energy basis of $H$ to $\tilde A_a$ enforcing detailed
  balance w.r.t.\ $\rho$.
  \item \textbf{Initial state $\sigma_0=\id/n$} is \emph{not} the proposal. It is
  only the starting point of the Markov chain $\sigma_{k+1}=\mathcal{M}[\sigma_k]$.
  The CPTP map $\mathcal{M}$ is built from $\{\tilde A_a\}$ and the reject Kraus
  $K=(\id-\sum_a\tilde A_a^\dagger\tilde A_a)^{1/2}$; it does not depend on
  $\sigma_0$. Stationarity is $\mathcal{M}[\rho]=\rho$; $\sigma_0$ affects only
  transient length (steps to convergence).
\end{itemize}

\subsection{Channel and convergence}
\begin{equation}
\mathcal{M}[\sigma]=\sum_a \tilde A_a\sigma\tilde A_a^\dagger
+ K\sigma K^\dagger .
\end{equation}
\texttt{run\_until\_converged} applies $\mathcal{M}$ until
$\dist(\sigma,\rho)\le\theta=\eps/(2R)$ or \texttt{gibbs\_max\_steps} is reached.
Per-step trace renormalization stabilizes stiff $M$. Optional cutoffs
$\{100,200,500\}$ snapshot $\sigma$ on the same trajectory for experiment~2.1.3.

% ===========================================================================
\section{Implementation choices}
\label{sec:impl}
% ===========================================================================

\begin{description}
  \item[CVXPY trace bound.] Default $R$ from~\eqref{eq:Rcvxpy}; certified
  $\Tr(Y^\star)\le R$.
  \item[Vectorized Bohr weights.] $w(\nu)=\tfrac12(1-\tanh(\nu/4))$ applied to
  the full Bohr matrix by array ufuncs (no per-element \texttt{vectorize}).
  \item[Parallel benchmark.] \texttt{run\_test.py} uses a process pool; each
  worker pins BLAS to one thread; worker count hard-capped at 64 on shared hosts.
  \item[Progress.] \texttt{tqdm} bars for calibration and oracle runs.
\end{description}

% ===========================================================================
\section{Benchmark: what we measure}
\label{sec:measure}
% ===========================================================================

\subsection{Calibration}
Bisection on $g$ (exact Gibbs) so one feasibility run takes
$\sim$\texttt{--target-iters} iterations (default 1000). No binary search at
run time.

\subsection{Experiments}
Epsilon sweep: $\eps\in\{0.1,\,0.05,\,0.01\}$; reference $\eps=0.01$ for 2.1.3
and 2.1.4.

\begin{table}[h]
\centering
\begin{tabular}{@{}llll@{}}
\toprule
Exp. & Violation & $\eps$ & Quantity \\
\midrule
2.1.1 & max & $0.1,0.05,0.01$ & Gibbs steps $k_t$ per oracle iter (spike plot) \\
2.1.2 & max & same & Wall-time pie (construction / channel / check / \ldots) \\
2.1.3 & max & $0.01$ & $(\eps/R)\dist(\sigma^{(c)},\sigma^{\mathrm{full}})$ at cutoffs \\
2.1.4 & random & $0.01$ & Class bars: avg max/mean steps, avg oracle iters \\
\bottomrule
\end{tabular}
\caption{Benchmark outputs under \texttt{results/<instance>/} and
\texttt{results/\_classes/}.}
\end{table}

\paragraph{Reading capped runs.}
If $\theta$ is still unreachable within \texttt{gibbs\_max\_steps}, spikes
flatline at the cap (lower bound on mixing time). After CVXPY $R$, this is less
common on instances where $\Tr(Y^\star)\ll 2|\textsc{opt}|$; large-trace
instances (e.g.\ \texttt{ThetaPrimeER}) may still require high caps or coarser
$\eps$.

% ===========================================================================
\section{Parameters and pitfalls}
\label{sec:params}
% ===========================================================================

\begin{description}
  \item[$R$.] Prefer~\eqref{eq:Rcvxpy}. Re-convert after changing
  \texttt{--R-margin}; delete old \texttt{results/} before re-benchmarking.
  \item[$\theta=\eps/(2R)$.] Scales mixing difficulty as $1/R^2$ at fixed $\eps$.
  \item[$g$.] Calibrates oracle \emph{iteration} count, not Gibbs steps per iter.
  \item[Channel cap.] \texttt{--gibbs-max-steps} is the main runtime knob; must
  exceed typical $k_t$ for informative 2.1.1 plots.
\end{description}

% ===========================================================================
\begin{thebibliography}{9}
\bibitem{vAG2019}
J.~van Apeldoorn and A.~Gily\'en,
\emph{Improvements in Quantum SDP-Solving with Applications},
arXiv:1804.05058.

\bibitem{GCDK2024}
A.~Gily\'en, C.-F.~Chen, J.~F.~Doriguello and M.~J.~Kastoryano,
\emph{Quantum generalizations of Glauber and Metropolis dynamics},
arXiv:2405.20322.
\end{thebibliography}

\end{document}
